\documentclass[a4paper]{article}

\usepackage[fontset=fandol]{ctex}
\usepackage{amsmath, amssymb}
\usepackage{graphicx}
\usepackage[margin=2.5cm]{geometry}
\usepackage[most]{tcolorbox}
\usepackage{listings}
\usepackage{hyperref}
\usepackage{float}

\newtcolorbox{knowledgebox}[1]{
    enhanced, colback=blue!5!white, colframe=blue!75!black, colbacktitle=blue!75!black,
    coltitle=white, fonttitle=\bfseries, title=#1, attach boxed title to top left={yshift=-2mm, xshift=2mm},
    boxrule=1pt, sharp corners
}

\newtcolorbox{warningbox}[1]{
    enhanced, colback=red!5!white, colframe=red!75!black, colbacktitle=red!75!black,
    coltitle=white, fonttitle=\bfseries, title=#1, sharp corners
}

\lstset{
    language=Python,
    basicstyle=\ttfamily\small,
    breaklines=true,
    frame=single,
    numbers=left,
    numberstyle=\tiny\color{gray},
    captionpos=b
}

\newcommand{\notetitle}{CS231n Lecture 11: Large Scale Distributed Training}
\newcommand{\noteauthors}{Codex Harness}
\newcommand{\notedate}{2026-04-15}
\newcommand{\videourl}{https://www.youtube.com/watch?v=9MvD-XsowsE}
\newcommand{\videocoverpath}{cover.jpg}

\begin{document}

\begin{titlepage}
\centering
{\Large 课程笔记\par}
\vspace{1.2cm}
{\huge\bfseries \notetitle\par}
\vspace{0.8cm}
{\large \noteauthors\par}
\vspace{0.3cm}
{\large \notedate\par}
\vspace{1.2cm}
\includegraphics[width=0.82\textwidth,height=0.45\textheight,keepaspectratio]{\videocoverpath}\par
\end{titlepage}

\tableofcontents
\newpage

\section{为什么大模型训练必然走向分布式}
% source_units: cu_11_01, cu_11_02
这讲的出发点非常实际：以 Llama3-405B 这类超大模型为例，单张 GPU 无论在显存、吞吐还是通信带宽上都已经无法独立承担训练任务。课程因此把一个 cluster 重新解释成“一台由很多 GPU 和高速互联组成的大计算机”。

这个视角很重要，因为它改变了训练问题的表述方式。我们不再问“怎么在一张卡上跑起来”，而是问：
\begin{itemize}
\item 数据应当切在哪里？
\item 模型参数应当切在哪里？
\item 序列与时间应当切在哪里？
\item 通信与计算应当怎样重叠？
\end{itemize}

\subsection{本章小结}
大规模训练的本质，不是简单地“多加几张卡”，而是把整个训练系统重新设计成一套分布式程序。

\section{数据并行: 最自然但不够用的第一步}
% source_units: cu_11_03, cu_11_04
数据并行（Data Parallelism）是最自然的起点：每张 GPU 持有一份完整模型副本，但各自处理不同 mini-batch 切片。

如果总 batch 大小为 $B$，每张卡上的局部梯度记作 $g^{(n)}$，那么聚合后的全局梯度可以写成：
\[
L(\theta) = \frac{1}{B}\sum_{i=1}^{B}\ell(x_i; \theta), \qquad
g = \frac{1}{N}\sum_{n=1}^{N} g^{(n)}
\]

这里的关键操作就是 \textbf{all-reduce}：每张卡先独立完成 forward/backward，再把梯度求和或求均值，最后同步更新参数。

\begin{figure}[H]
\centering
\includegraphics[width=0.86\textwidth]{pdf_pages/page-32.png}
\caption{数据并行把样本分到多卡，但每张卡仍保留完整模型副本。}
\end{figure}

\begin{lstlisting}[caption={数据并行训练 step 的最小伪代码}]
for local_batch in loader_shard:
    loss = model(local_batch).loss()
    loss.backward()
    all_reduce_gradients(model.parameters())
    optimizer.step()
    optimizer.zero_grad()
\end{lstlisting}

它的优点是概念简单、易于扩展；但它也很快触碰上限。原因很直接：\textbf{模型参数、梯度、优化器状态都被完整复制在每张卡上。} 当模型足够大时，仅复制开销就会压垮显存。

\subsection{本章小结}
数据并行解决的是算力扩展问题，但没有解决模型状态复制问题。

\section{FSDP 与显存优化: 把状态切开保存}
% source_units: cu_11_05, cu_11_06
当数据并行不再够用时，下一步自然是对模型状态做分片。课程在 slides 中出现了 ``FSPD'' 的写法，但其实际内容对应的是 \textbf{Fully Sharded Data Parallelism, FSDP}。

FSDP 的核心思想是：不要让每张卡都完整保存参数、梯度和优化器状态，而是把这些状态切开存。这样每张卡只负责自己那一片，在需要计算某层时再做 gather，在计算结束后再释放或重新切分。

\begin{figure}[H]
\centering
\includegraphics[width=0.86\textwidth]{pdf_pages/page-46.png}
\caption{FSDP 把参数、梯度和优化器状态切分到不同设备，从复制问题转化为通信与调度问题。}
\end{figure}

这类方法带来的收益是显著的：
\begin{itemize}
\item 更大的模型可以塞进同样数量的 GPU。
\item 可用 batch size 往往也随之提升。
\item 为后续与 pipeline/tensor/context parallel 组合打下基础。
\end{itemize}

课程还提醒，很多看似次要的细节也会影响显存：例如最后一层激活是否保留、某些张量生命周期是否拉得过长、是否采用 hybrid sharding 等。分布式训练并不是“选一种 parallelism 就结束”，而是大量低层工程细节的总和。

\subsection{本章小结}
FSDP 的关键不是“更高级”，而是它首次真正解决了状态复制的根本问题。

\section{Activation Checkpointing, HFU 与 MFU}
% source_units: cu_11_07, cu_11_08
即使做了 sharding，激活（activations）本身仍可能成为瓶颈。于是课程引出 \textbf{activation checkpointing}：训练时不保存所有中间激活，而是在反向传播需要时重新计算一部分 forward。

\begin{figure}[H]
\centering
\includegraphics[width=0.86\textwidth]{pdf_pages/page-75.png}
\caption{Activation checkpointing 用额外计算换显存，是大模型训练中的常规工程手段。}
\end{figure}

这是一种很典型的 tradeoff：
\begin{itemize}
\item 少存激活，显存压力下降。
\item 反向阶段需要重算，计算量上升。
\end{itemize}

课程随后区分了两个吞吐指标：
\[
\mathrm{HFU} = \frac{t_{\text{theoretical matmul}}}{t_{\text{actual}}}, \qquad
\mathrm{MFU} = \frac{t_{\text{theoretical model}}}{t_{\text{actual}}}
\]

直观上，HFU 更接近“硬件上的矩阵乘法利用率”，而 MFU 更接近“整个模型训练系统实际有多高效”。课程之所以强调 MFU，是因为 checkpointing、helper kernels、通信等待等开销，往往不会被 HFU 充分反映。

\begin{warningbox}{为什么不能只看 HFU}
某个系统可能 matmul 很忙，看起来 HFU 不差，但如果通信阻塞、重算太多或 pipeline 气泡很大，整体训练速度依然不理想。MFU 更接近端到端系统表现。
\end{warningbox}

\subsection{本章小结}
checkpointing 解决的是“显存怎么省”，HFU/MFU 回答的是“训练有没有真的跑快”。两者共同决定大规模训练的工程质量。

\section{Context, Pipeline, Tensor: 把模型沿不同维度切开}
% source_units: cu_11_09, cu_11_10, cu_11_11
当模型、序列和层数都继续增长时，只靠数据并行和 sharding 仍然不够。课程于是把并行策略扩展到了更多维度。

\textbf{Context Parallelism} 主要针对长上下文 Transformer。它把序列切开，让不同 GPU 负责不同上下文片段，从而缓解长序列注意力的显存压力。

\begin{figure}[H]
\centering
\includegraphics[width=0.86\textwidth]{pdf_pages/page-121.png}
\caption{Context parallel 与 pipeline parallel 表明，模型不仅能沿数据维切，也能沿序列维和层维切。}
\end{figure}

\textbf{Pipeline Parallelism} 则把模型不同层分配到不同 GPU，上游卡处理完 microbatch 后立刻把中间激活送给下游卡。课程特别强调了 \textbf{pipeline bubble}：如果 microbatch 不足或调度不佳，某些设备会空转，最大 MFU 会被流水线深度限制。

\textbf{Tensor Parallelism} 更进一步，它不切样本也不切层，而是直接切单层内部的大矩阵乘法。以 Transformer 为例，projection 和 MLP 中的大矩阵往往就是最适合切分的对象。

\begin{figure}[H]
\centering

\includegraphics[width=0.86\textwidth]{pdf_pages/page-135.png}
\caption{Tensor parallel 直接把单层内部的大矩阵运算拆到多卡，是扩展超大宽度模型的重要方式。}
\end{figure}

课程最后用 \textbf{ND Parallelism} 收束全讲：真正的大模型训练几乎总是多种并行策略的组合，而不是只选其中一种。数据维、参数维、层维、序列维和算子维都可能同时参与切分。

\subsection{本章小结}
并行不是单选题，而是多维组合题。真正的系统设计能力，体现在能否把这些维度拼成一个高效而稳定的整体。

\section{总结与延伸}
这讲的中心思想可以浓缩为一句话：\textbf{大规模训练，本质上是把训练程序沿多个维度切开，并用通信把它重新粘起来。}

从数据并行开始，我们看到完整模型复制的上限；从 FSDP 和 checkpointing，我们看到显存如何通过切分和重算换回来；从 context / pipeline / tensor parallel，我们看到模型、序列和层本身都可以成为切分轴。

对后续课程与实践而言，最重要的不是记住每个缩写，而是建立一个系统观：
\begin{itemize}
\item 先找出真正的瓶颈是显存、通信还是吞吐。
\item 再决定应沿数据、状态、层、序列或算子哪个维度切分。
\item 最后用 MFU 这类更接近端到端效率的指标衡量整体方案。
\end{itemize}

\end{document}
